\chapter{Obsah elektronické přílohy}

Elektronická příloha obsahuje zdrojové texty práce, projekty hardwaru
v~programu KiCad, konfigurace firmwaru zařízení a logy z~ladění. Aktuální
verze všech souborů je dostupná také v~repozitáři na GitHubu:

\makeatletter
\g@addto@macro\UrlBreaks{\do\_}
\makeatother
\begin{center}
  \small\url{https://github.com/cryptonhd2004/bap-smart_home_thread_esp}
\end{center}

\section*{Použité nástroje a verze}

Hardware byl navržen v~programu KiCad ve verzi~10.0. Firmware zařízení byl
vytvořen v~nástroji ESPHome [DOPLNIT: verze] s~externí komponentou
\texttt{esphome-matter}~\cite{esphome-matter}, která do ESPHome přidává
podporu protokolu Matter. Komponenta je postavena na sadě
ESP-Matter~\cite{espressif-esp-matter} a podporuje pouze čipy ESP32
s~frameworkem ESP-IDF, který ESPHome spravuje automaticky. Použitá verze
komponenty je určena commitem [DOPLNIT: hash] repozitáře. Přenos dat
v~síti Thread zajišťuje komponenta \texttt{openthread}
nástroje ESPHome~\cite{esphome-openthread}.

Zařízení jsou zprovozněna pomocí souborů ve formátu YAML. Čidlo teploty
pracuje jako plnohodnotné zařízení Thread (FTD), zatímco vypínač pracuje
v~minimálním režimu (MTD). Podklady pro výrobu desek plošných spojů jsou
součástí adresářů jednotlivých zařízení.

\bigskip

{\small
\dirtree{%.
.1 /\DTcomment{kořenový adresář přiloženého archivu}.
.2 README.md\DTcomment{popis projektu a stav práce}.
.2 bap\_text\DTcomment{zdrojové soubory práce (\LaTeX)}.
.3 sablona-prace.tex\DTcomment{hlavní soubor pro sazbu práce}.
.3 nastaveni.tex\DTcomment{údaje o práci, abstrakt, klíčová slova}.
.3 text\DTcomment{jednotlivé kapitoly a seznam literatury}.
.3 obrazky\DTcomment{obrázky použité v~textu}.
.2 chytre\_svetlo\DTcomment{projekt KiCad: chytré světlo (schéma, DPS)}.
.2 Chytre\_Tlacitko\_BMS\DTcomment{projekt KiCad: bateriový vypínač (schéma, DPS)}.
.2 esp\_thread\DTcomment{projekt KiCad: chytrý teploměr}.
.2 test\_builds\DTcomment{konfigurace ESPHome a testovací firmware}.
.3 test-teplomer.yaml\DTcomment{konfigurace teploměru a porovnávaných čidel}.
.3 tlacitko-thread.yaml\DTcomment{konfigurace vypínače}.
.3 c6-matter-teplomer-firmware.factory.bin\DTcomment{binární firmware teploměru}.
.2 logs\_for\_debuging\DTcomment{logy z~ladění}.
}
}